\documentclass{article}
\usepackage{graphicx} % Required for inserting images

\title{Introducción}
\author{Daniel Gago}
\date{January 2026}

\begin{document}

\maketitle

\section{Los primeros SBC}
    \subsection{Dendral}
        El primer modelo que usa conocimiento experto.

    \subsection{MYCIN}
        Este modelo se caractriza por separar el conocimiento de resolución del problema(control) del conocimiento del dominio.\\

        Ademas de esto, este modelo usa un esquema de razonamiento bajo incertidumbre.\\

    \subsection{PROSPECTOR}
        Este modelo tambien usa un esquema bajo incertidumbre, auqnue diferente al de MYCIN, ademas que tuvo un gran exito comercial.

\section{Características de los Sistemas Inteligentes}

    Los Sistemas Inteligentes (SI) se diferencian de los sistemas informáticos clásicos en que están diseñados para abordar dominios y problemas de mayor complejidad, donde la información puede ser incompleta, imprecisa o cambiante, y no siempre existe una solución determinista claramente definida.
        
    \subsection{Dominios y problemas complejos}
        
    Los SI se aplican a problemas que superan las capacidades de la informática sistemática tradicional. Estos dominios suelen requerir razonamiento, adaptación y el uso de heurísticas, como ocurre en áreas tales como el diagnóstico, la planificación, los sistemas expertos o la toma de decisiones.
        
    \subsubsection{Sistemas procedimentales frente a Sistemas Basados en Conocimiento}
        
    Los sistemas informáticos clásicos son sistemáticos y procedimentales: reciben una entrada de datos organizada y producen una salida mediante algoritmos deterministas, de modo que un mismo conjunto de datos de entrada genera siempre el mismo resultado.
        
    En contraste, los Sistemas Basados en Conocimiento (SBC) adoptan un enfoque declarativo y heurístico. Estos sistemas se adaptan al estado del problema, pueden modificar sus objetivos durante la ejecución y emplean mecanismos de razonamiento para optimizar su comportamiento.
        
    \subsection{Mantenimiento perfectivo continuo}
        
    Una característica esencial de los Sistemas Inteligentes es la posibilidad de realizar un mantenimiento perfectivo continuo. Esto implica que el sistema puede evolucionar progresivamente mediante la incorporación de nuevo conocimiento o la mejora de heurísticas, sin necesidad de rediseñar completamente el software.
        
    \subsection{Separación entre conocimiento y métodos de resolución}
        
    Los SI se basan en la separación entre el conocimiento específico del dominio y los métodos de resolución del problema. Esta separación favorece la reutilización del software, ya que un mismo motor de inferencia puede aplicarse a distintos dominios simplemente cambiando la base de conocimiento.
        
    \subsection{Justificación y explicación}
        
    Otra característica fundamental es la capacidad de justificación y explicación. Los Sistemas Inteligentes deben ser capaces de explicar las decisiones tomadas y las conclusiones alcanzadas, proporcionando transparencia y aumentando la confianza del usuario en el sistema.
        
    \subsection{Conocimiento explícito e implícito}
        
    En los Sistemas Inteligentes se distingue entre conocimiento explícito, que se representa directamente mediante hechos, reglas u otras estructuras formales, y conocimiento implícito, que está incorporado en el comportamiento del sistema o en sus algoritmos. Un buen diseño de SI busca maximizar la representación explícita del conocimiento para facilitar su comprensión y mantenimiento.

\section{Ventajas y limitaciones de los Sistemas Basados en Conocimiento}

    Los Sistemas Basados en Conocimiento (SBC) presentan una serie de ventajas que los hacen especialmente adecuados para la resolución de problemas complejos. No obstante, también poseen limitaciones inherentes a su enfoque y a las técnicas empleadas.

    \subsection{Ventajas}

    Entre las principales ventajas de los Sistemas Basados en Conocimiento se encuentran las siguientes:

    \begin{itemize}
        \item \textbf{Mantenimiento del conocimiento}: el conocimiento se representa de forma explícita, lo que facilita su actualización, corrección y ampliación.
        \item \textbf{Resolución de problemas complejos}: permiten abordar dominios mal definidos o con gran complejidad mediante razonamiento y heurísticas.
        \item \textbf{Ajuste de objetivos}: el sistema puede modificar sus objetivos o estrategias durante la ejecución en función del estado del problema.
        \item \textbf{Tratamiento de la incertidumbre}: incorporan mecanismos para manejar información incompleta o imprecisa, como factores de certeza o razonamiento probabilístico.
        \item \textbf{Explicación del razonamiento}: son capaces de justificar las decisiones tomadas, aumentando la transparencia y la confianza del usuario.
        \item \textbf{Reducción de costes}: permiten automatizar tareas que normalmente requerirían expertos humanos, reduciendo costes a largo plazo.
        \item \textbf{Aumento de la fiabilidad}: ofrecen soluciones consistentes y evitan errores derivados del cansancio o la variabilidad humana.
        \item \textbf{Modularidad}: la separación entre conocimiento y mecanismos de inferencia favorece el desarrollo modular del sistema.
    \end{itemize}
    
    \subsection{Limitaciones}
    
    A pesar de sus ventajas, los SBC presentan también una serie de limitaciones:
    
    \begin{itemize}
        \item \textbf{Dificultad en la adquisición del conocimiento}: extraer y formalizar el conocimiento de expertos humanos es un proceso complejo y costoso.
        \item \textbf{Reutilización del conocimiento}: el conocimiento suele estar muy ligado a un dominio concreto, lo que dificulta su reutilización directa en otros contextos.
        \item \textbf{Falta de creatividad y sentido común}: los SBC se limitan al conocimiento previamente representado y no poseen capacidades creativas ni de razonamiento de sentido común.
        \item \textbf{Obstáculos para el aprendizaje y la adaptación}: muchos SBC tradicionales no aprenden automáticamente y requieren intervención humana para su actualización.
    \end{itemize}


\section{Metodologías}
    Con todo esto, para ayudar a que los modelos de IC fueran mejores y evitar problemas se empezaron a usar aproximaciones a las metodologías usadas por la IS. Sin embargo estas metodologías de la IS no tenían en cuenta el almacenamiento o obtencion del conocimiento que precisan los modelos de IC. \\

    Con todo esto empiezan a generarse metodologías para aplicarse a la IC, sin embargo estas metodologías son muy precisas y complejas. Ademas de que tras su aplicación podemos decir que antes el problema es el conocimiento, pero ahora el problema es otro, el sentido, pues quien decide lo que es etico o donde esta la linea entre la buena aplicación o mala aplicación de los modelos de esta disciplina.

\section{Separación del conocimiento/representación}
    Para poder afrontar los problemas con el conocimiento en la IC lo que se hizo fue separar el conocimiento en un nuevo nivel, lo que ayuda a mejorar la comprension y entendimiento que tenemos sobre el. \\

    Tal como decia Newell, con mirar solo el codigo no basta para entender lo que hace,hay que describir qué sabe y qué objetivos persigue, de ahi CommonKADS convierte las ideas de Newell en una metodología. \\

    Estas metodologías son importantes ya que permiten explicar, reutilizar y validar con expertos. \\

\section{MLR}
    Este metodo se olvida del conocimiento del problema y representa de manera genérica el conocimiento de resolución de problemas. 

\section{Tareas}
    Tipo de problema o conjunto de instancias de problemas con algo en común\\

    \subsection{Métodos y Subtareas}
        Medios que nos permiten realizar una tarea.\\

        Método: Conjunto de subtareas que permiten transformar el estado inicial en el estado objetivo de una tarea.
    \subsection{Estructura de tareas}
        Árbol de tareas, métodos y subtareas aplicadas recursivamente hasta que las tareas que se alcancen puedan llevarse a cabo mediante el conocimiento disponible.
    

    

    
    

    
    
        

\end{document}
